%=====================================================================
\section{K3 lattice realisation of the $140$-label representation}\label{sec:k3}
%=====================================================================
This section uses only the label representation of Theorem~\ref{thm:F} at $s=7$; nothing in it concerns $H_s$. \emph{Category~A (cited)} means proved
conditional only on the published theorems named, which we use as stated in the cited sources and do not re-prove.

\paragraph{Notation.} Throughout, $A=\left(\begin{smallmatrix}2&7\\7&7\end{smallmatrix}\right)$ is the Gram matrix of $Q_7$ (\S\ref{sec:delta}), $B(u,v)=u^{\mathsf T}Av$, and a lattice is written with its \emph{norm} $x\cdot x$ (so $Q_2=\tfrac12x\cdot x$). We write $N_1$, $N_2$ (italic) for two even lattices.

%---------------------------------------------------------------------
\subsection{The even lattice $N_2$ and its discriminant form}\label{sec:disc}
The form $(\Z^2,A)$ is odd ($Q_7\notin\Z$), so it is not itself a K3 lattice. Its even sublattice of index $2$ is
\[
N_2:=\{x\in\Z^2:\ x_2\equiv0\bmod2\},\qquad \mathrm{Gram}\big((1,0),(0,2)\big)=\begin{pmatrix}2&14\\14&28\end{pmatrix},\qquad\det=-140 .
\]
The other natural even version is $N_1:=(\Z^2,2A)$, with Gram matrix $\left(\begin{smallmatrix}4&14\\14&14\end{smallmatrix}\right)$, also of determinant $-140$.

\begin{proposition}[Discriminant form of $N_2$]\label{prop:disc8}
\begin{enumerate}[label=(\roman*)]
\item In the basis $(1,0),(-7,2)$ of $N_2$ (base change $P=\left(\begin{smallmatrix}1&-7\\0&1\end{smallmatrix}\right)$ from $(1,0),(0,2)$), $N_2\cong\langle2\rangle\oplus\langle-70\rangle$, i.e.\ $x\cdot x=2(x_1+7y)^2-70y^2$. In particular $N_2$ is even, of signature $(1,1)$.
\item $A_{N_2}=N_2'/N_2\cong\Z/2\times\Z/70$, $|A_{N_2}|=140$, generated by $(\tfrac12,0)$ and $(0,\tfrac1{70})$ in the $\langle2\rangle\oplus\langle-70\rangle$ basis, with $q_{N_2}(x,y)=\tfrac{x^2}4-\tfrac{y^2}{140}\bmod1$.
\item The local Jordan components and their Gauss-sum signatures are:
\begin{center}\small
\begin{tabular}{llll}\toprule
$p$ & group & $q$ on the $p$-part & signature mod $8$\\\midrule
$2$ & $(\Z/2)^2$, generators $(\tfrac12,0),(0,\tfrac12)$ & $\langle\tfrac14\rangle\oplus\langle\tfrac14\rangle$, $q=(x^2+y^2)/4$ & $2$\\
$5$ & $\Z/5$ & $\langle\tfrac35\rangle\cong\langle\tfrac25\rangle$ & $4$\\
$7$ & $\Z/7$ & $\langle\tfrac27\rangle\cong\langle\tfrac17\rangle$ & $2$\\\bottomrule
\end{tabular}
\end{center}
The local signatures sum to $2+4+2\equiv0=\mathrm{sign}(N_2)\bmod8$ (oddity formula). The $5$- and $7$-parts are exactly the local forms $q_5=2x^2/5$ and $q_7=x^2/7$ of Theorem~\ref{thm:tensor7}.
\end{enumerate}
\end{proposition}
\begin{proof}
(i) $(1,0)\cdot(1,0)=2$, $(-7,2)\cdot(-7,2)=98-196+28=-70$, $(1,0)\cdot(-7,2)=-14+14=0$, and $\det P=1$. (ii) is read off from (i): $\langle2\rangle'=\tfrac12\Z$ with $q=\tfrac12\cdot2\cdot\tfrac14$, $\langle-70\rangle'=\tfrac1{70}\Z$ with $q=-\tfrac1{140}$. (iii) The $2$-part is generated by $(\tfrac12,0)$ and $(0,\tfrac{35}{70})$, with $q(0,\tfrac12)=-35^2/140\equiv\tfrac14$; the $5$-part by $(0,\tfrac15)$, $q=-14^2/140\equiv\tfrac35$; the $7$-part by $(0,\tfrac17)$, $q=-10^2/140\equiv\tfrac27$. The isomorphisms hold because $3\cdot2^{-1}\equiv4$ is a square mod $5$ and $2\equiv3^2$ mod $7$. The signatures are the standard Gauss-sum values.
\end{proof}
\catA{Elementary; exact.}

\subsection{The 140-label representation is the Weil representation of $N_2$}\label{sec:weil}
\paragraph{Conventions.} For an even lattice $L$ we use Borcherds' convention \cite{Borcherds}: $\rho_L(T)\e_\gamma=\e(q(\gamma))\e_\gamma$ and $\rho_L(S)\e_\gamma=\e(-\mathrm{sign}(L)/8)\,|A_L|^{-1/2}\sum_\delta\e(-b(\gamma,\delta))\e_\delta$; the dual $\rho_L^*$ is the complex conjugate. Since $\mathrm{sign}(N_2)=0$, $\rho_{N_2}$ is an honest representation of $\mathrm{SL}_2(\Z)$ and $\e(\mathrm{sign}/8)=1$. On the label side, the labels are $(\alpha,\beta)$ with $\alpha\in(L'\cup(a+L'))/\Z^2$ ($70$ classes), $\beta\in\{0,\beta_1\}$, $a=(\tfrac7{10},-\tfrac15)$ and $\beta_1$ the stored label of $-a$ mod $\Z^2$; the operators $S_{\rm lab}$, $T_{\rm lab}$ are built in the convention of Remark~\ref{rem:conv} (with exact rational phases, all denominators dividing $140$). Following Remark~\ref{rem:Sconv}, the label representation in standard normalisation is
\[
\rho_{\rm lab}:\quad S\mapsto S_{\rm lab}^{-1},\qquad T\mapsto T_{\rm lab}.
\]

\begin{theorem}[Weil realisation of the label representation]\label{thm:weil8}
\begin{enumerate}[label=(\roman*)]
\item $\rho_{\rm lab}\cong\rho_{N_2}$ as representations of $\mathrm{SL}_2(\Z)$. No character twist is needed and no metaplectic sign arises.
\item An explicit intertwiner is the following. Index the basis of $\rho_{N_2}$ by $z\in N_2'/N_2$, written in $\Z^2$-coordinates, and let $[z]=z\bmod\Z^2$. Every class $\alpha$ has exactly two lifts $z$, $z+(0,1)$ in $N_2'/N_2$. Put
\[
X\,\e_z\;=\;u_{([z],0)}\;+\;\e\big(B(z,\beta_1)\big)\,u_{([z],\beta_1)} .
\]
Then $X\rho_{N_2}(S)=S_{\rm lab}^{-1}X$, $X\rho_{N_2}(T)=T_{\rm lab}X$, and $X$ is invertible.
\item $\rho_{\rm lab}\not\cong\rho_{N_2}^*$, $\rho_{\rm lab}\not\cong\rho_{N_1}$ and $\rho_{\rm lab}\not\cong\rho_{N_1}^*$.
\end{enumerate}
\end{theorem}
\begin{proof}
(ii) All entries of $\sqrt{140}\,X\rho_{N_2}(S)$, $\sqrt{140}\,S_{\rm lab}^{-1}X$ (note $\sqrt{140}=2\sqrt{35}$ and $\sqrt{35}\,S_{\rm lab}^{-1}$ has entries $\e(\cdot)$), $X\rho_{N_2}(T)$ and $T_{\rm lab}X$ lie in $\Z[\zeta_{280}]$. Each of the $2\cdot140^2$ entry identities is the vanishing of an explicit integer combination of powers of $\zeta_{280}$; it was checked by exact reduction modulo the cyclotomic polynomial $\Phi_{280}$, with $0$ failures for $S$ and $0$ for $T$. $X$ is block diagonal with respect to $\alpha$; the block of the two lifts $z,z'=z+(0,1)$ is $\left(\begin{smallmatrix}1&1\\\e(B(z,\beta_1))&\e(B(z',\beta_1))\end{smallmatrix}\right)$, with determinant $\e(B(z,\beta_1))\big(\e(B((0,1),\beta_1))-1\big)\neq0$ because $B((0,1),\beta_1)=7(-\tfrac7{10})+7\cdot\tfrac15\equiv\tfrac12\bmod1$. (i) follows from (ii); $X$ intertwines the untwisted $\rho_{N_2}$, and $\mathrm{sign}(N_2)=0$.
(iii) $T_{\rm lab}$ is monomial; its label permutation has $70$ fixed points and $35$ two-cycles, and its eigenvalue multiset ($36$ distinct values) equals that of $\rho_{N_2}(T)$ but differs (exactly, as multisets of elements of $\Q/\Z$) from those of $\rho_{N_2}^*(T)$, $\rho_{N_1}(T)$ and $\rho_{N_1}^*(T)$. Conjugate representations have conjugate $T$-spectra.
\end{proof}
\catA{Exact computation in $\Z[\zeta_{280}]$ ($0$ failures among the $2\cdot140^2$ entry identities; $B(v,\beta_1)\equiv\tfrac12$; $X$ invertible) and exact
$T$-spectrum comparison. A second exact solution, with offsets $\{0,\tfrac45\}$, is another element of the Hom space of Proposition~\ref{prop:hom8}. Further $T$-spectrum
tests: $2\langle5,-7\rangle$, $2\langle-5,7\rangle$, $2\left(\begin{smallmatrix}3&1\\1&-12\end{smallmatrix}\right)$ and its negative do not match; $2\langle1,-35\rangle$
and $2\left(\begin{smallmatrix}1&1\\1&-34\end{smallmatrix}\right)$ match only because both are isometric to $N_2$, and their negatives match the dual.}
\begin{proposition}[Relations and Hom spaces modulo $\ell$]\label{prop:hom8}
Let $\ell=840841\equiv1\bmod280$. Computing modulo $\ell$ (with a fixed primitive $280$-th root of unity in $\mathbb F_\ell$):
\begin{enumerate}[label=(\roman*)]
\item $(S_{\rm lab}^{-1},T_{\rm lab})$ satisfies $S^4=I$ and $(ST)^3=S^2$, as $\rho_{N_2}$ does; $(S_{\rm lab},T_{\rm lab})$ satisfies $S^4=I$ but \emph{fails} $(ST)^3=S^2$, so it is not a representation of $\mathrm{SL}_2(\Z)$ in this normalisation (confirming Remark~\ref{rem:Sconv}).
\item $\dim_{\mathbb F_\ell}\mathrm{Hom}(\rho_{N_2},(S_{\rm lab}^{-1},T_{\rm lab}))=8$, and a random element has rank $140$.
\item $\dim_{\mathbb F_\ell}\mathrm{Hom}(\rho_{N_2},(S_{\rm lab},T_{\rm lab}))=4$, and a random element has rank $72$: the two normalisations agree only on the $72$-dimensional subspace $W_P$.
\end{enumerate}
\end{proposition}
\catA{Exact computation modulo $\ell$. The isomorphism itself is Theorem~\ref{thm:weil8}. By semicontinuity the characteristic-$0$ Hom dimensions are at most $8$
and $4$. \emph{Category~B (reading):} the value $8$ is consistent with eight pairwise non-isomorphic irreducible sectors $W_5^\pm\otimes W_7^\pm\otimes\{A_3,B_1\}$
(\S\ref{sec:family}); that the characteristic-$0$ dimension equals $8$ is not separately proved here.}
\begin{remark}[Dual, twist, metaplectic sign]\label{rem:twist8}
$\rho_{\rm lab}\not\cong\rho_{N_2}^*$ (Theorem~\ref{thm:weil8}(iii)); no twist by a character of $\mathrm{SL}_2(\Z)$ is needed; $\mathrm{sign}(N_2)=0$, so there is no metaplectic sign. The Weil representation depends only on the discriminant form, so any even lattice with discriminant form $\cong q_{N_2}$ and signature $\equiv0\bmod8$ (e.g.\ $N_2\oplus U$, $N_2\oplus E_8$, and the lattice $T$ of Proposition~\ref{prop:T8}, up to dualisation) gives the same representation.
\end{remark}

\subsection{The transcendental lattice}\label{sec:T}
Let $\Lambda_{K3}=U^3\oplus E_8(-1)^2$ (even unimodular, signature $(3,19)$), and let $U\oplus U\subset\Lambda_{K3}$ have bases $e_1,f_1,e_2,f_2$ with $e_i\cdot f_i=1$ and all other products $0$.
\begin{proposition}[Explicit primitive embedding and $T$]\label{prop:T8}
Put $n_1=e_1+f_1$ (norm $2$), $n_2=e_2-35f_2$ (norm $-70$), $t_1=e_1-f_1$ (norm $-2$), $t_2=e_2+35f_2$ (norm $70$).
\begin{enumerate}[label=(\roman*)]
\item $\mathrm{span}(n_1,n_2)\cong N_2$ and is primitive in $U\oplus U$ (the $2\times2$ minors of $(n_1\ n_2)$ have gcd $1$).
\item $t_1,t_2$ are orthogonal to $n_1,n_2$ and span the full orthogonal complement of $\mathrm{span}(n_1,n_2)$ in $U\oplus U$ (index $1$; their Gram determinant is $-140$).
\item Hence the orthogonal complement of $N_2$ in $\Lambda_{K3}$ is
\[
T:=N_2^\perp=\langle-2\rangle\oplus\langle70\rangle\oplus U\oplus E_8(-1)^2,\qquad\mathrm{sign}(T)=(2,18),\qquad A_T\cong\Z/2\times\Z/70,
\]
with $q_T(x,y)=-\tfrac{x^2}4+\tfrac{y^2}{140}=-q_{N_2}(x,y)$ under the identity identification of generators. The gluing is the one induced by the unimodular overlattice $U\oplus U$ (consistent with Nikulin \cite{Nikulin}, Prop.~1.6.1).
\end{enumerate}
\end{proposition}
\begin{proof}
Direct computation: $n_1\cdot n_1=2$, $n_2\cdot n_2=-70$, $n_1\cdot n_2=0$, $t_1\cdot t_1=-2$, $t_2\cdot t_2=70$, and $n_i\cdot t_j=0$. The complement of a rank-$2$ primitive sublattice of the unimodular $U\oplus U$ has rank $2$ and $|\det|=140$; $\mathrm{span}(t_1,t_2)$ lies in it with $|\det|=140$, so the index is $1$. The other summands $U\oplus E_8(-1)^2$ are orthogonal to $U\oplus U$. The discriminant form of $\langle-2\rangle\oplus\langle70\rangle$ is read off as in Proposition~\ref{prop:disc8}.
\end{proof}
\catA{Elementary; exact. Nikulin's existence theorem is not needed, since the embedding is explicit.}
\begin{corollary}[The residual law in K3 terms]\label{cor:T8}
$\rho_{\rm lab}\;\cong\;\rho_{N_2}\;\cong\;\rho_T^*=\overline{\rho_T}$ (Borcherds' convention; honest $\mathrm{SL}_2(\Z)$-representations).
\end{corollary}
\begin{proof}
$\mathrm{sign}(T)=-16\equiv0\bmod8$ and $q_T=-q_{N_2}$ (hence $b_T=-b_{N_2}$) under the identification of Proposition~\ref{prop:T8}, so in the formulas for $\rho_T(S)$ and $\rho_T(T)$ every phase is the conjugate of the corresponding phase for $\rho_{N_2}$. Combine with Theorem~\ref{thm:weil8}.
\end{proof}
\catA{Purely lattice-theoretic; no cited geometric theorem is used. In words: the $140$-label residual law transforms under the dual Weil representation of a signature-$(2,18)$ lattice, which (Theorem~\ref{thm:K3exist8}) is the transcendental lattice of an $N_2$-polarised K3 surface. This is a statement about a shared \emph{representation}; it is not a statement that any residual series is a K3 generating function (see Remark~\ref{rem:open8}).}

\subsection{Geometric realisation}\label{sec:geom}
\begin{theorem}[$N_2$-polarised K3 surfaces]\label{thm:K3exist8}
\begin{enumerate}[label=(\roman*)]
\item The primitive embedding $N_2\hookrightarrow\Lambda_{K3}$ is unique up to $O(\Lambda_{K3})$; hence $T$ is unique up to isometry.
\item There is a projective K3 surface $X$ with $\mathrm{NS}(X)\cong N_2$ and $T_X\cong T$.
\item $N_2$ represents neither $-2$ nor $0$ (non-trivially). Hence such an $X$ contains no smooth rational curves and admits no elliptic fibration, and its ample cone is the whole component of the positive cone containing an ample class.
\end{enumerate}
\end{theorem}
\begin{proof}[Proof sketch]
(i) Nikulin \cite[Thm.~1.14.4]{Nikulin} (p.~126 of the translation): $\Lambda_{K3}$ is even unimodular of signature $(3,19)$, $N_2$ is even of signature $(1,1)$, $3>1$, $19>1$, and $\mathrm{rk}\,\Lambda_{K3}-\mathrm{rk}\,N_2=20\ge2+l(A_{N_2})=4$. (ii) Morrison \cite[Cor.~2.9(i)]{Morrison} (p.~112; proved there via Cor.~1.9): for $\rho\le10$ every even lattice of signature $(1,\rho-1)$ is the N\'eron--Severi lattice of some algebraic K3 surface (equivalently, surjectivity of the period map); with (i), $T_X\cong T$. (iii) $x\cdot x=2\big((x_1+7y)^2-35y^2\big)$; $-2$ would need $u^2-35y^2=-1$, impossible since $-1$ is not a square mod $7$; $0$ would need $35$ to be a square. Without $(-2)$-classes the ample cone is a full component of the positive cone, and an elliptic fibration requires a non-zero isotropic class in $\mathrm{NS}$ (standard, \cite{Huybrechts}).
\end{proof}
\catA{\textbf{Category~A (cited)}: we use Nikulin \cite[Thm.~1.14.4]{Nikulin} and Morrison \cite[Cor.~2.9(i)]{Morrison}, and standard facts on ample cones and
elliptic fibrations \cite{Huybrechts}. The arithmetic in (iii) is exact.}

\subsection{Automorphisms}\label{sec:aut}
Let $\gamma=\left(\begin{smallmatrix}-1&-7\\2&13\end{smallmatrix}\right)$, the Pell automorph of Theorem~\ref{thm:B}: $\gamma^{\mathsf T}A\gamma=A$, $\det\gamma=1$, $\mathrm{tr}\,\gamma=12$, eigenvalues $6\pm\sqrt{35}$. It preserves $N_2$ (the second coordinate of $\gamma x$ is $2x_1+13x_2$). Let $r=\mathrm{diag}(1,-1)$ in the $\langle2\rangle\oplus\langle-70\rangle$ basis; there $\gamma=\left(\begin{smallmatrix}6&35\\1&6\end{smallmatrix}\right)$ (multiplication by $\varepsilon=6+\sqrt{35}$).

\begin{proposition}[Action on $A_{N_2}$ and $O(N_2)$]\label{prop:gamma8}
\begin{enumerate}[label=(\roman*)]
\item The induced map $\bar\gamma$ on $A_{N_2}$ acts on the local parts as follows: on the $2$-part $(\Z/2)^2$ as a non-trivial involution (not the identity, which on this part equals $\pm1$); on the $5$-part as $+1$; on the $7$-part as $-1$. Hence $\bar\gamma\neq\pm1$ and $\bar\gamma^2=\mathrm{id}$. This matches the action of $\gamma$ in Theorem~\ref{thm:tensorodd}(ii) on the $7$/$5$/$2$ parts.
\item $O(N_2)=\{\pm\gamma^k,\ \pm r\gamma^k:\ k\in\Z\}$, and $r\gamma r=\gamma^{-1}$.
\item On $A_{N_2}$: $\gamma^{2j}$ acts as $+1$ and $r\gamma^{2j}$ as $-1$, both preserving the positive cone; odd powers $\pm\gamma^{2j+1}$, $\pm r\gamma^{2j+1}$ act as neither $+1$ nor $-1$; $-\gamma^{k}$ and $-r\gamma^k$ do not preserve the positive cone.
\end{enumerate}
\end{proposition}
\begin{proof}
(i), (iii): exact computation on the $140$ classes. (ii): in the diagonal basis $x\cdot x=2(u^2-35y^2)$. Proper automorphs of $u^2-35y^2$ correspond to solutions of $t^2-140w^2=4$, i.e.\ to norm-$1$ units of $\Z[\sqrt{35}]$, which are $\pm\varepsilon^k$ with $\varepsilon=6+\sqrt{35}$ fundamental; there is no unit of norm $-1$ ($u^2-35y^2=-1$ is impossible mod $7$). Improper automorphs are $r$ times proper ones, and $r\gamma r=\left(\begin{smallmatrix}6&-35\\-1&6\end{smallmatrix}\right)=\gamma^{-1}$.
\end{proof}
\catA{Exact.}
\begin{theorem}[Automorphisms of $N_2$-polarised K3 surfaces]\label{thm:aut8}
Let $X$ be as in Theorem~\ref{thm:K3exist8}, with $\mathrm{NS}(X)=N_2$ glued to $T_X=T$ as in Proposition~\ref{prop:T8}.
\begin{enumerate}[label=(\roman*)]
\item $\gamma^2\oplus\mathrm{id}_{T}$ extends to a Hodge isometry of $H^2(X,\Z)$ preserving the ample cone, and is induced by a unique $f\in\mathrm{Aut}(X)$. $f$ has infinite order; $f^*$ has spectral radius $(6+\sqrt{35})^2=71+12\sqrt{35}$ on $H^2(X,\Z)$ (a quadratic unit, not a Salem number of higher degree), so its topological entropy is
\[
h(f)=2\log(6+\sqrt{35})=4\sqrt{35}\,I\approx4.9558,
\]
where $I=\log(6+\sqrt{35})/(2\sqrt{35})$ is the constant of Theorem~\ref{thm:C}.
\item $r\oplus(-\mathrm{id}_T)$ is induced by a non-symplectic involution $\iota\in\mathrm{Aut}(X)$.
\item If $\mathrm{NS}(X)\cong N_2$ (so that the Hodge isometries of $T_X$ are only $\pm1$, Proposition~\ref{prop:liftgamma}(i)), then $\mathrm{Aut}(X)=\langle f,\iota\rangle\cong\Z\rtimes\Z/2$ (the infinite dihedral group $D_\infty$), and $\gamma$ itself does not lift to $\mathrm{Aut}(X)$.
\end{enumerate}
\end{theorem}
\begin{proof}[Proof sketch]
By Nikulin's gluing criterion \cite{Nikulin}, Cor.~1.5.2, an isometry $g_{\rm NS}\oplus g_T$ of $N_2\oplus T$ extends to $H^2(X,\Z)$ iff $\bar g_{\rm NS}$ and $\bar g_T$ agree under the glue isomorphism $A_{N_2}\cong A_T$. By Proposition~\ref{prop:gamma8}, $\bar\gamma^2=+1$ and $\bar r=-1$, compatible with $\mathrm{id}_T$ and $-\mathrm{id}_T$, which are Hodge isometries. Both preserve the positive cone, hence (Theorem~\ref{thm:K3exist8}(iii)) the ample cone. By the strong Torelli theorem (Piatetski-Shapiro--Shafarevich, Burns--Rapoport; e.g.\ \cite{Huybrechts}, Thm.~7.5.3) they are induced by unique automorphisms $f$, $\iota$; $\iota$ acts by $-1$ on $H^{2,0}$, so it is non-symplectic. The eigenvalues of $f^*$ are $(6\pm\sqrt{35})^2$ on $\mathrm{NS}\otimes\R$ and $1$ on $T$, and the entropy is the logarithm of the spectral radius (Gromov--Yomdin \cite{GY}). For (iii): every automorphism acts on $T_X$ by $\pm1$, so its restriction to $N_2$ preserves the positive cone and acts on $A_{N_2}$ by $\pm1$; by Proposition~\ref{prop:gamma8}(iii) these restrictions are exactly $\gamma^{2j}$ and $r\gamma^{2j}$, which form $\langle\gamma^2\rangle\rtimes\langle r\rangle\cong D_\infty$; the action on $H^2$ is injective by Torelli. Since $\bar\gamma\neq\pm1$, $\gamma$ does not lift.
\end{proof}
\catA{\textbf{Category~A (cited)}: we use the strong Torelli theorem, Nikulin's gluing criterion \cite[Cor.~1.5.2]{Nikulin}, the Gromov--Yomdin entropy formula, and
Proposition~\ref{prop:liftgamma}(i) for (iii). The lattice inputs (Proposition~\ref{prop:gamma8}) are exact. Numerically
$2\log(6+\sqrt{35})=2\log(11.91608\ldots)=4.9558\ldots$. The identity $2\log(6+\sqrt{35})=4\sqrt{35}\,I$ is definitional; \emph{Category~B (reading):} that the
regulator driving the Baker obstruction is ``the same'' as this entropy is an observation about definitions, and no theorem linking $C$ to $X$ is claimed.}

\begin{proposition}[Lifting $\gamma$]\label{prop:liftgamma}
\begin{enumerate}[label=(\roman*)]
\item If $\mathrm{NS}(X)\cong N_2$, then the Hodge isometries of $T_X$ are $\pm1$, and $\gamma$ does not lift to $\mathrm{Aut}(X)$.
\item If $\gamma$ lifts on some $X$ with $N_2\subset\mathrm{NS}(X)$ primitive and $M=N_2^\perp\subset\mathrm{NS}(X)$, then $M$ carries the $2$- and $7$-parts of $A_{N_2}$ with $\mathrm{rk}\,M\equiv4\pmod8$, or the $2$- and $5$-parts with $\mathrm{rk}\,M\equiv6\pmod8$. Hence $\rho(X)\ge6$.
\item Let $M=A_1(-1)^2\oplus K_7$ with $K_7=\left(\begin{smallmatrix}-2&1\\1&-4\end{smallmatrix}\right)$. There is an even rank-$6$ gluing $N_2\oplus M\subset L$ and an isometry $g_M$ of $M$ such that $\gamma\oplus g_M\oplus\mathrm{id}_{T}$ is a Hodge isometry of $H^2(X,\Z)$ for every K3 surface $X$ with $\mathrm{NS}(X)\cong L$ (such $X$ exist by \cite[Cor.~2.9(i)]{Morrison}). On this rank-$6$ Noether--Lefschetz locus some $w\circ(\gamma\oplus g_M)$, $w\in W(L)$, is induced by an automorphism, so $\mathrm{Aut}$ grows there.
\end{enumerate}
\end{proposition}
\begin{proof}[Proof sketch]
(i) The reduction $\bar\gamma\in O(A_{N_2})$ has order $2$: it swaps on the $2$-part, is $+1$ on the $5$-part and $-1$ on the $7$-part, so $\bar\gamma\neq\pm1$. The Hodge isometries of $T_X$ form a cyclic group $\mu_m$, with $T_X\otimes\Q$ a $\Q(\zeta_m)$-vector space. If $m>2$, a prime $p\mid m$ gives, after passing to the kernel of the action on $A_T$, an element of order $p$ without fixed vectors acting trivially on $A_T$. This forces $pA_T=0$, which is impossible because $A_T\cong\Z/2\times\Z/70$ is not $p$-elementary. So every automorphism acts on $T_X$ by $\pm1$, hence on $A_{N_2}$ by $\pm1$, and $\gamma$ cannot lift.
(ii) The glue must carry $\bar\gamma$ to an isometry of $A_M$. Milgram's formula, with Gauss-sum signatures $4,2,4,6,6\bmod8$ for the relevant $p$-parts, gives the two congruences.
(iii) An exact lattice computation gives the anti-isometric glue and $g_M$ (the swap on $A_1(-1)^2$ and $-1$ on $K_7$) with $\bar g_M\circ\phi=\phi\circ\bar\gamma$. Because $L$ contains $(-2)$-vectors, $\gamma\oplus g_M\oplus\mathrm{id}$ need not preserve the ample chamber. The strong Torelli theorem \cite[Thm.~7.5.3]{Huybrechts} supplies a Weyl element $w$ for which $w\circ(\gamma\oplus g_M\oplus\mathrm{id})$ is induced by an automorphism.
\end{proof}
\catA{(i), (ii): written proofs. (iii): exact lattice computation, together with \cite[Cor.~2.9(i)]{Morrison} and the strong Torelli theorem (cited), so it is conditional on these citations. \emph{Open:} whether $w=1$ can be taken, i.e.\ whether $\gamma$ itself lifts on this locus.}

\begin{remark}[Category~B and open items for this section]\label{rem:open8}
None of the following is a result of this paper.
\begin{enumerate}[label=(\alph*)]
\item \emph{Noether--Lefschetz / Heegner series} (Category~B). The Noether--Lefschetz/Heegner generating series of the $N_2$-polarised family are vector-valued modular forms of weight $1+18/2=10$ for $\rho_T$ or $\rho_T^*$ (depending on convention; Borcherds, Kudla--Millson, Maulik--Pandharipande, cited by name only). The residual objects of Theorems~\ref{thm:F}--\ref{thm:G} are weight-$0$ Maass-type objects with $\lambda=\tfrac14$. A shared \emph{representation} (Corollary~\ref{cor:T8}) is not a shared \emph{form}; no theta-lift or Borcherds-product link has been checked.
\item \emph{Special members} (open). Whether $\gamma$ itself lifts on the rank-$6$ locus of Proposition~\ref{prop:liftgamma}(iii), or on any $X$ with $\rho(X)\ge6$.
\item \emph{Interpretation} (Category~B). Any physical or ``unification'' reading of $X$, $f$ or $\iota$.
\item \emph{Scope.} Every K3 statement in this paper concerns $N_2$, because the $140$-label representation is $\rho_{N_2}$ and not $\rho_{N_1}$
(Theorem~\ref{thm:weil8}(iii)).
\end{enumerate}
\end{remark}
